\section{Alexandrov Phase Nerves and Finite Realization}
\label{app:phase-nerve}

This research addendum concerns arbitrary complete FDE value profiles on
relational frames. It does not impose probabilistic thresholds or persistence
of the positive and negative valuation channels. It is separate from the
probability-path results in the main text.

\paragraph{Definitions and convention.}
Let $Q=\{\Tval,\Fval,\Bval,\Nval\}$ and let $v:W\to Q$.
For a reflexive, transitive relation $R$, use the successor Alexandrov topology,
so that $\operatorname{Int}(A)=\{w: \forall u\,(wRu\Rightarrow u\in A)\}$.
Write $X_q=v^{-1}(q)$ and
\[
 S_w=\{q\in Q: \exists u\,(wRu\ \wedge\ v(u)=q)\}.
\]
The witness-based phase nerve is
\[
 \mathcal N_{R,v}
 =\{\sigma\subseteq Q: \exists w\in W\;
          \forall q\in\sigma\; w\in\operatorname{Cl}(X_q)\}.
\]
When $W$ is nonempty, this includes the empty face. A phase whose fibre is empty
does not appear as a vertex. No claim about the homotopy type of $W$ follows
merely from this definition.

\begin{theorem}[Successor-profile classification]
For every $\sigma\subseteq Q$,
\[
 \sigma\in\mathcal N_{R,v}
 \quad\Longleftrightarrow\quad
 \exists w\in W\;\sigma\subseteq S_w.
\]
\end{theorem}
\begin{proof}
Dualizing successor interior gives
$w\in\operatorname{Cl}(X_q)$ if and only if some successor of $w$ has value $q$.
Substitution into the definition of the nerve gives the equivalence.
\end{proof}

\begin{theorem}[Universal finite equivalence-frame realization]
Let $K\subseteq\mathcal P(Q)$ be downward closed and have at least one vertex.
There exists a finite reflexive, symmetric, transitive frame and an FDE profile
whose phase nerve equals $K$, including its empty face and absent vertices.
\end{theorem}
\begin{proof}
Use incidence worlds
\[
 W_K=\{(\tau,q):\tau\in K,\ q\in\tau\},\qquad
 v(\tau,q)=q.
\]
Put $(\tau,q)R(\rho,r)$ precisely when $\tau=\rho$.
This is an equivalence relation. A world labelled by $\tau$ sees exactly the
phases in $\tau$. Consequently,
\[
 \sigma\in\mathcal N_{R,v}
 \quad\Longleftrightarrow\quad
 \exists\tau\in K\;(\tau\ne\varnothing\ \wedge\ \sigma\subseteq\tau).
\]
Downward closure gives the forward implication to $\sigma\in K$.
For a nonempty $\sigma\in K$, choose $\tau=\sigma$; for the empty face, use
the assumed vertex. This proves equality, not just containment of nerves.
\end{proof}

\paragraph{Finite implementation and trust boundary.}
The Lean construction encodes faces by four Boolean flags and enumerates all
face/phase candidates. The formal bound is a list of at most $64$ worlds
containing every world. Selection from an arbitrary proposition-valued complex
uses classical decidability; this is not a claim of an executable decision
procedure for an unspecified input predicate. A focused axiom report is emitted
when the realization module is built. The new proofs introduce no project axiom,
\texttt{sorry}, or native-evaluation proof step.

\paragraph{Boundary and significance.}
The complex consisting only of the empty face is excluded: any world witnesses
the singleton containing its own value. On an empty world type the witness-based
nerve has no faces at all. Thus the vertex hypothesis cannot simply be replaced
by the assertion that $K$ has a face.

The construction shows that even equivalence-frame (S5) accessibility imposes
no restriction on the abstract phase-complex shapes, when world count and
valuation are allowed to vary. This does not say that every complex can be
realized on one fixed frame. Restrictions may arise from fixed accessibility,
valuation regularity, persistence, or formula-specific threshold geometry.
Nor is the realizing frame claimed to be $T_0$ or antisymmetric. The ingredients
are standard; priority for this particular combination remains unaudited.

\paragraph{Formal correspondence.}
The classification theorem is
\path{alexandrov_phase_simplex_iff}. Its module is:
\par\noindent{\small\path{PEL4.TopologicalEvidenceAlexandrovNerve}}\par

The realization module is:
\par\noindent{\small\path{PEL4.TopologicalEvidenceNerveRealization}}\par
Its principal declarations are:
\begin{itemize}
\item \path{every_phase_complex_realized};
\item \path{phaseRealizationWorlds_complete};
\item \path{phaseRealizationWorlds_length_le};
\item \path{phaseRealization_symmetric}.
\end{itemize}
